\documentclass[11pt,a4paper]{article}

% -------------------------------------------------
% Paquetes
% -------------------------------------------------
\usepackage[utf8]{inputenc}
\usepackage[T1]{fontenc}
\usepackage[spanish]{babel}
\usepackage{geometry}
\usepackage{enumitem}
\usepackage{hyperref}
\usepackage{xcolor}
\usepackage{listings}
\usepackage{booktabs}
\usepackage{array}
\usepackage{tabularx}

\geometry{
  top=2.5cm,
  bottom=2.5cm,
  left=2.7cm,
  right=2.7cm
}

\hypersetup{
  colorlinks=true,
  linkcolor=black,
  urlcolor=blue
}

\setlist[itemize]{noitemsep, topsep=4pt}
\setlist[enumerate]{noitemsep, topsep=4pt}

\lstset{
  basicstyle=\ttfamily\small,
  frame=single,
  breaklines=true,
  columns=fullflexible
}

\lstset{
  literate=
  {á}{{\'a}}1 {é}{{\'e}}1 {í}{{\'i}}1 {ó}{{\'o}}1 {ú}{{\'u}}1
  {Á}{{\'A}}1 {É}{{\'E}}1 {Í}{{\'I}}1 {Ó}{{\'O}}1 {Ú}{{\'U}}1
  {ñ}{{\~n}}1 {Ñ}{{\~N}}1 {¿}{{?`}}1 {¡}{{!`}}1
}

% -------------------------------------------------
% Documento
% -------------------------------------------------
\title{
  \textbf{Protocolo de Trabajo con Git, GitHub y Linear}\\
  \large Taller de Integración II y Taller de Integración IV
}
\author{}
\date{Segundo semestre 2026}

\begin{document}

\maketitle

\section{Objetivo}

Esta guía fija la forma en que los equipos de Taller de Integración II (TI2) y
Taller de Integración IV (TI4) trabajarán en el proyecto.

El flujo busca dos cosas concretas: que el trabajo se pueda seguir desde Linear
hasta GitHub y que cada equipo conserve sus responsabilidades técnicas.

El proyecto usa tres herramientas:

\begin{itemize}
  \item \textbf{Linear}: planificación, asignación y seguimiento del trabajo.
  \item \textbf{Git}: control de versiones.
  \item \textbf{GitHub}: alojamiento del código, Pull Requests, revisión de
    código e integración continua.
\end{itemize}

La división de responsabilidades es directa:

\begin{center}
  \textbf{Linear gestiona el trabajo y GitHub gestiona el código.}
\end{center}

% -------------------------------------------------

\section{Organización de los equipos}

El proyecto estará compuesto por dos equipos.

\subsection{Taller de Integración II}

TI2 tendrá seis estudiantes y se encargará de:

\begin{itemize}
  \item desarrollo de la aplicación web;
  \item desarrollo del backend y API mediante arquitectura por capas;
  \item implementación de los casos de uso y reglas de negocio compartidas;
  \item diseño e implementación de la base de datos;
  \item pruebas asociadas a sus componentes;
  \item despliegue de sus propios servicios;
  \item documentación técnica correspondiente.
\end{itemize}

\subsection{Taller de Integración IV}

TI4 tendrá tres estudiantes y se encargará de:

\begin{itemize}
  \item desarrollo de la aplicación móvil;
  \item planificación general del proyecto;
  \item preparación y priorización inicial del trabajo;
  \item coordinación entre ambos equipos;
  \item control de calidad;
  \item revisión de código;
  \item seguimiento de dependencias;
  \item apoyo técnico e inducción al equipo TI2 cuando sea necesario.
\end{itemize}

TI4 debe ayudar a TI2 a resolver sus tareas cada vez con más autonomía. En una
revisión, TI4 señala el problema y explica el criterio. La persona autora hace
el cambio en su propia rama.

% -------------------------------------------------

\section{Repositorio GitHub}

Todo el producto vivirá en \textbf{un único repositorio GitHub}. El monorepo
contendrá las aplicaciones, el backend y los paquetes compartidos.

\begin{lstlisting}
Proyecto
|
+-- .github/workflows/       Integración continua
+-- apps/
|   +-- web/                 Aplicación web de TI2
|   +-- mobile/              Aplicación móvil de TI4
|
+-- convex/                  Backend por capas y base de datos compartidos
|   +-- presentation/        Funciones públicas de Convex
|   +-- application/         Casos de uso
|   +-- domain/              Reglas de negocio
|   +-- infrastructure/      Persistencia y servicios externos
+-- packages/                Código compartido cuando corresponda
+-- package.json
+-- bun.lock
\end{lstlisting}

Cada equipo desarrolla, prueba y despliega los componentes que tiene asignados
dentro del monorepo.

Un repositorio común no elimina la separación de responsabilidades. TI2
mantiene la aplicación web y el backend. TI4 mantiene la aplicación móvil.

La integración entre los equipos pasa por la API pública de la capa de
Presentación del backend. Ambas aplicaciones consumen sus tipos generados. No
se copian implementaciones ni se mantienen contratos externos.

\subsection{Arquitectura por capas}

Las dependencias deben avanzar en esta dirección:

\begin{lstlisting}
Presentación
      |
      v
Aplicación
      |
      v
Dominio
      ^
      |
Infraestructura
\end{lstlisting}

En el backend, las funciones públicas de Convex reciben la solicitud y llaman
a un caso de uso. Los casos de uso coordinan la operación. El Dominio contiene
las reglas. Infraestructura conecta con la persistencia, la autenticación y los
servicios externos.

Una branch o Pull Request que cruce varias capas debe explicar qué cambia en
cada una. React, Mobile y las funciones públicas no deben resolver necesidades
de negocio con consultas directas a la base de datos ni con funciones extensas.

% -------------------------------------------------

\section{Uso de Linear}

Todo el proyecto usará \textbf{un único workspace de Linear}.

Dentro de este workspace existirán dos Teams:

\begin{itemize}
  \item \textbf{TI2}: Web, Backend y Base de Datos.
  \item \textbf{TI4}: Gestión y Aplicación Móvil.
\end{itemize}

Los dos Teams forman parte del mismo proyecto.

\subsection{Triage}

TI4 hará el \textit{triage} inicial. Para cada issue debe:

\begin{itemize}
  \item crear o revisar los issues necesarios;
  \item asignarlos al Team correspondiente;
  \item definir su prioridad;
  \item verificar que tengan una descripción suficiente;
  \item establecer criterios de aceptación;
  \item identificar dependencias entre ambos equipos;
  \item asignarlos al Cycle correspondiente.
\end{itemize}

Cuando un issue esté listo en el backlog, el equipo responsable organiza su
ejecución.

\subsection{Linear como fuente de verdad}

Toda tarea técnica que requiera trabajo significativo debe tener un issue en
Linear.

WhatsApp, Discord y otros canales sirven para conversar, no para asignar trabajo
de manera formal. Si una conversación genera una tarea, alguien debe registrarla
en Linear.

\subsection{Sub-issues}

El proyecto no usa \textbf{sub-issues} de Linear.

Para una tarea con varios pasos pequeños, usa una lista en la descripción del
issue.

Por ejemplo:

\begin{lstlisting}
TI2-37 Implementar autenticación

- Crear endpoint de login
- Validar credenciales
- Generar token
- Agregar pruebas
- Documentar endpoint
\end{lstlisting}

Si una parte del trabajo necesita asignación, seguimiento o revisión propios,
conviértela en un issue independiente.

Así se evita gastar issues del plan en tareas que no necesitan seguimiento
separado.

% -------------------------------------------------

\section{Cycles y Sprints}

Cada \textit{Cycle} de Linear representa un Sprint. No crearemos un Cycle por
semana.

La correspondencia es esta:

\begin{lstlisting}
Cycle 1 -> Sprint 1
Cycle 2 -> Sprint 2
Cycle 3 -> Sprint 3
\end{lstlisting}

Las semanas de cada Sprint quedan para seguimiento, reuniones y control de
avance.

Al cerrar un Cycle, el equipo revisa:

\begin{itemize}
  \item trabajo completado;
  \item trabajo pendiente;
  \item issues bloqueados;
  \item deuda técnica generada;
  \item dependencias entre TI2 y TI4;
  \item trabajo que debe trasladarse al siguiente Sprint.
\end{itemize}

% -------------------------------------------------

\section{Estados de los Issues}

El workflow tendrá pocos estados:

\begin{center}
  \texttt{Backlog}
  $\rightarrow$
  \texttt{Todo}
  $\rightarrow$
  \texttt{In Progress}
  $\rightarrow$
  \texttt{In Review}
  $\rightarrow$
  \texttt{Done}
\end{center}

También puedes usar:

\begin{center}
  \texttt{Canceled}
\end{center}

Los estados significan lo siguiente:

\begin{description}
  \item[Backlog:] trabajo identificado pero todavía no comprometido.
  \item[Todo:] trabajo seleccionado para el Cycle actual.
  \item[In Progress:] existe trabajo activo sobre el issue.
  \item[In Review:] existe uno o más Pull Requests pendientes de revisión.
  \item[Done:] todo el trabajo requerido por el issue fue integrado.
  \item[Canceled:] el trabajo dejó de ser necesario.
\end{description}

Cuando sea posible, la integración Linear--GitHub cambiará estos estados de
forma automática.

% -------------------------------------------------

\section{Integración de Linear con GitHub}

El repositorio del monorepo debe conectarse al workspace de Linear.

Cada issue conserva el Team responsable, TI2 o TI4. El identificador distingue
el área aunque todo el código viva en el mismo repositorio.

Los identificadores de Linear relacionan el código con el trabajo.

Por ejemplo:

\begin{lstlisting}
TI2-37
TI4-18
\end{lstlisting}

Un Pull Request puede referenciar el issue desde la branch, el título o la
descripción.

La trazabilidad queda así:

\begin{center}
  \textbf{
    Issue Linear
    $\rightarrow$
    Branch
    $\rightarrow$
    Commits
    $\rightarrow$
    Pull Request
    $\rightarrow$
    Review
    $\rightarrow$
    Merge
  }
\end{center}

% -------------------------------------------------

\section{Estrategia de ramas}

El repositorio usará \textbf{feature branches pequeñas y de corta duración}.

La única rama permanente es:

\begin{lstlisting}
main
\end{lstlisting}

No habrá una rama permanente \texttt{develop}. Tampoco usaremos Git Flow con
ramas permanentes de desarrollo, release o hotfix.

\subsection{Creación de una rama}

Antes de empezar una tarea:

\begin{enumerate}
  \item confirma que existe un issue en Linear;
  \item actualiza la rama \texttt{main};
  \item crea una rama desde \texttt{main};
  \item trabaja solo en el objetivo del issue.
\end{enumerate}

Flujo conceptual:

\begin{lstlisting}
main
 |
 +-- rama TI2-37-login-web
 |
 +-- rama TI2-42-auth
 |
 +-- rama TI4-18-login-mobile
\end{lstlisting}

% -------------------------------------------------

\section{Nombres de ramas}

Usa, de preferencia, el nombre de branch que genera Linear.

El nombre debe conservar:

\begin{itemize}
  \item nombre o identificador del desarrollador;
  \item identificador del issue de Linear;
  \item una descripción breve.
\end{itemize}

Ejemplo:

\begin{lstlisting}
vrivera/ti2-37-login
\end{lstlisting}

Si Linear genera un nombre demasiado largo, puedes abreviarlo.

Por ejemplo, en lugar de:

\begin{lstlisting}
vrivera/ti2-37-implementar-endpoint-de-autenticacion-de-usuarios
\end{lstlisting}

usa:

\begin{lstlisting}
vrivera/ti2-37-login
\end{lstlisting}

Conserva siempre el identificador de Linear.

No uses:

\begin{lstlisting}
login
feature1
cambios
prueba
rama-nueva
\end{lstlisting}

% -------------------------------------------------

\section{Tamaño de las ramas}

Una rama debe resolver \textbf{una intención concreta}.

Mantén las ramas pequeñas para facilitar:

\begin{itemize}
  \item revisión;
  \item pruebas;
  \item resolución de conflictos;
  \item integración continua;
  \item detección de errores;
  \item reversión de cambios.
\end{itemize}

No fijaremos un límite artificial de líneas modificadas. El criterio es que otra
persona pueda entender y revisar el Pull Request sin encontrarse con varias
funcionalidades independientes.

% -------------------------------------------------

\section{Pull Requests}

Todo cambio que llegue a \texttt{main} debe pasar por un \textbf{Pull Request}
en el repositorio del monorepo.

Nadie puede hacer push directo a \texttt{main}.

\subsection{Título del Pull Request}

El título debe seguir la convención del repositorio. Revisa los Pull Requests
recientes y el historial de Git antes de elegirlo. Con \textbf{Squash and
Merge}, el título suele convertirse en el mensaje del commit que llega a
\texttt{main}; por eso debe ser breve, claro y explicar el resultado del cambio.

Incluye el identificador de Linear cuando corresponda y describe qué mejora
obtiene el usuario o el sistema. No enumeres archivos, funciones o pasos de
implementación en el título.

Evita:

\begin{lstlisting}
TI2-37: Agregar AuthForm, query login, mutation login, schema y pruebas
\end{lstlisting}

Prefiere:

\begin{lstlisting}
TI2-37: Permitir que los usuarios inicien sesión desde la Web
\end{lstlisting}

\subsection{Descripción del Pull Request}

Abre la descripción con una explicación sencilla del problema que originó el
issue. Después explica la solución y añade los detalles necesarios para
revisarla. No empieces con un inventario de archivos, funciones eliminadas o
componentes modificados.

Evita:

\begin{lstlisting}
Se agregó AuthForm, se creó la mutation login, se modificó schema.ts,
se actualizaron los hooks y se añadieron pruebas para cada componente.
\end{lstlisting}

Prefiere:

\begin{lstlisting}
El formulario de acceso no permitía iniciar sesión desde la Web. Ahora la
aplicación valida las credenciales mediante el backend compartido y muestra
los errores de autenticación sin duplicar la lógica del servidor.
\end{lstlisting}

El Pull Request debe indicar:

\begin{itemize}
  \item qué problema resuelve;
  \item qué cambios introduce;
  \item cómo puede probarse;
  \item issue de Linear relacionado;
  \item cualquier consideración relevante para el reviewer.
\end{itemize}

Ejemplo:

\begin{lstlisting}
TI2-37: Implementar login de usuarios
\end{lstlisting}

La descripción debe dar el contexto suficiente para revisar el cambio sin
reconstruirlo desde el código.

% -------------------------------------------------

\section{Relación entre Issues y Pull Requests}

Un issue y un Pull Request no tienen que mantener una relación 1:1.

Un issue representa una \textbf{unidad de trabajo}.

Un Pull Request representa una \textbf{unidad de integración de código}.

Un issue puede requerir varios Pull Requests.

Ejemplo:

\begin{lstlisting}
TI2-42 Autenticación

    |
    +-- PR 1: modelo de usuario
    |
    +-- PR 2: endpoint login
    |
    +-- PR 3: refresh token
\end{lstlisting}

El issue termina cuando todo el trabajo necesario está integrado.

% -------------------------------------------------

\section{Stacked Pull Requests}

Si no puedes dividir una tarea en cambios independientes, encadena Pull
Requests o usa \textit{stacked PRs}.

Ejemplo:

\begin{lstlisting}
main
 |
 +-- PR A
      |
      +-- PR B
           |
           +-- PR C
\end{lstlisting}

Esta técnica es una excepción. No es el flujo predeterminado.

Como límite práctico:

\begin{center}
  \textbf{usa como máximo 3 stacked PRs; el límite absoluto es 4.}
\end{center}

Si una cadena supera cuatro Pull Requests dependientes, revisa la planificación
y divide la tarea de nuevo.

% -------------------------------------------------

\section{Revisión de código}

Otra persona debe revisar cada Pull Request antes de integrarlo.

Una revisión puede terminar en:

\begin{itemize}
  \item aprobación;
  \item comentarios;
  \item solicitud de cambios.
\end{itemize}

La revisión debe comprobar:

\begin{itemize}
  \item cumplimiento del issue;
  \item corrección;
  \item claridad del código;
  \item mantenibilidad;
  \item errores potenciales;
  \item pruebas;
  \item seguridad;
  \item respeto de la dirección de dependencias entre capas;
  \item ubicación correcta de las reglas de negocio;
  \item compatibilidad con otros componentes;
  \item contrato de API cuando corresponda.
\end{itemize}

% -------------------------------------------------

\section{Progresión de las revisiones de TI2}

TI2 está en una etapa inicial de formación, así que la supervisión cambiará
durante el semestre.

\subsection*{Etapa inicial}

Durante la primera etapa, integrantes de TI4 revisarán los Pull Requests de TI2.

TI4 puede pedir cambios, explicar el problema y sugerir una solución. La persona
autora hace los cambios en su propia rama.

\subsection*{Etapa intermedia}

En la etapa intermedia, TI2 empezará a revisar código entre pares.

TI4 seguirá revisando directamente los cambios relacionados con:

\begin{itemize}
  \item API;
  \item arquitectura;
  \item base de datos;
  \item seguridad;
  \item integración con móvil;
  \item cambios con impacto significativo.
\end{itemize}

\subsection*{Etapa avanzada}

En la etapa avanzada, TI2 debería revisar e integrar sus cambios con autonomía.

TI4 se concentrará en:

\begin{itemize}
  \item coordinación;
  \item QA;
  \item arquitectura;
  \item integración;
  \item resolución de dependencias.
\end{itemize}

% -------------------------------------------------

\section{Política de Merge}

Usaremos \textbf{Squash and Merge}. Así, cada Pull Request deja un solo commit
en \texttt{main}, aunque la branch tenga commits intermedios.

Después del merge:
\begin{itemize}
  \item conserva la rama de desarrollo;
  \item actualiza el issue relacionado;
  \item incluye la documentación asociada.
\end{itemize}

Las ramas integradas forman parte del historial de trabajo y de la evidencia de
participación de cada integrante.
% -------------------------------------------------

\section{Protección de \texttt{main}}

TI4 mantendrá protegida la rama \texttt{main} en GitHub. Ustedes no tienen que
configurar estas reglas cada vez que trabajen, pero sí deben conocerlas para
saber por qué todos los cambios pasan por Pull Request.

GitHub bloqueará como mínimo:

\begin{itemize}
  \item el push directo a \texttt{main};
  \item el merge sin Pull Request;
  \item el merge sin al menos una aprobación;
  \item el merge cuando fallen las validaciones automáticas requeridas;
  \item el merge con conversaciones pendientes.
\end{itemize}

Las mismas reglas aplican a TI2 y TI4. Incluso los cambios de TI4 pasan por
Pull Request para que el historial sea visible y ustedes puedan participar de
la revisión.

% -------------------------------------------------

\section{Integración Continua}

La integración continua es el filtro automático que les ayuda a detectar
errores antes del merge. Cada Pull Request activa GitHub Actions y ustedes
pueden revisar el resultado desde la pestaña de comprobaciones.

Las rutas modificadas determinan qué verificaciones se ejecutan. Como mínimo,
esperamos que la CI:

\begin{itemize}
  \item instale las dependencias desde el \texttt{package.json} y el único
    \texttt{bun.lock} de la raíz;
  \item ejecute el build de la aplicación o paquete afectado;
  \item aplique las herramientas de formato;
  \item ejecute el linter;
  \item ejecute las pruebas unitarias del Dominio y de los casos de uso;
  \item ejecute las pruebas de integración de Infraestructura y Presentación
    cuando corresponda.
\end{itemize}

Si su cambio afecta el backend, Web, Mobile o un paquete compartido, la CI
revisa todos los componentes afectados.

Si la CI falla, revisen el log, corrijan el problema en su branch y vuelvan a
publicar el cambio. No integren un Pull Request con la CI fallando salvo que TI4
acuerde una excepción justificada.

% -------------------------------------------------

\section{Integración entre TI2 y TI4}

Ustedes, como TI2, mantienen la aplicación Web y los servicios compartidos. El
detalle técnico de esas responsabilidades, la arquitectura y los adaptadores
está en \texttt{tech.tex}; aquí nos interesa cómo coordinar los cambios con TI4.

El monorepo permite que TI2 y TI4 coordinen una funcionalidad en el mismo Pull
Request o separen los cambios en varios Pull Requests. Elijan la opción que
facilite la revisión y el despliegue, y relacionen siempre los cambios en
Linear.

Cuando un cambio de ustedes afecte una funcionalidad de TI4:

\begin{itemize}
  \item relacionen el cambio con un issue de Linear;
  \item identifiquen al equipo afectado;
  \item comuniquen los cambios incompatibles antes de integrarlos;
  \item actualicen la documentación del servicio;
  \item indiquen si afecta Presentación, Aplicación, Dominio o Infraestructura;
  \item revisen la autenticación, los tipos y la estructura de datos cuando
    corresponda.
\end{itemize}

% -------------------------------------------------

\section{Cambios en el modelo de datos}

Traten cualquier cambio del modelo de datos como un cambio que puede afectar a
ambos equipos. La estructura técnica del modelo y su ubicación están explicadas
en \texttt{tech.tex}; en este flujo nos importa que el cambio sea visible y
coordinado.

Antes de integrarlo, registren el cambio en Linear, indiquen qué funcionalidades
puede afectar y avisen a TI4 si requiere adaptar Mobile.

% -------------------------------------------------

\section{Dependencias entre equipos}

Si su tarea depende del trabajo de TI4, o si TI4 queda esperando un cambio de
ustedes, registren la dependencia en Linear. Esto permite que TI4 la siga y que
ustedes sepan qué información o entrega falta.

Ejemplo:

\begin{lstlisting}
TI4-25 Pantalla de recuperación de clave

Bloqueada por:

TI2-61 Endpoint de recuperación de clave
\end{lstlisting}

TI4 hará seguimiento de estas dependencias durante la coordinación del proyecto
y les avisará si una de ellas bloquea el avance.

% -------------------------------------------------

\section{Definición de Done}

Para TI4, un issue no termina porque ``el código funciona en el computador del
desarrollador''. Termina cuando otra persona puede revisar, integrar y
entender el cambio dentro del proyecto.

Marquen una tarea como \texttt{Done} cuando:

\begin{itemize}
  \item cumpla sus criterios de aceptación;
  \item el código esté integrado mediante Pull Request;
  \item las observaciones de la revisión estén resueltas;
  \item la CI haya finalizado correctamente;
  \item hayan realizado las pruebas necesarias;
  \item hayan actualizado la documentación cuando corresponda;
  \item las dependencias entre capas respeten la arquitectura definida;
  \item no queden dependencias conocidas sin registrar;
  \item el cambio esté disponible en \texttt{main}.
\end{itemize}

% -------------------------------------------------

\section{Flujo de trabajo completo}

Este es el recorrido que esperamos que sigan para cada tarea:

\begin{enumerate}
  \item TI4 prepara el issue y lo asigna al Team y al Cycle correspondientes.
  \item Tomen la tarea desde Linear y revisen qué resultado se espera.
  \item Identifiquen las capas afectadas y el punto de entrada de la
    funcionalidad.
  \item Actualicen localmente \texttt{main} y creen desde Linear una branch
    asociada al issue dentro del monorepo.
  \item Implementen el cambio y hagan commits en su branch.
  \item Publiquen la branch en GitHub y abran el Pull Request.
  \item Relacionen el Pull Request con el issue en Linear.
  \item Revisen el resultado de la CI y respondan las observaciones del review.
  \item Cuando CI y la revisión estén aprobadas, TI4 hará Squash and Merge.
  \item Conserven la branch para mantener visible el historial de trabajo.
  \item Actualicen el estado del trabajo en Linear.
  \item Coordinen el despliegue cuando el cambio lo requiera.
\end{enumerate}

El flujo queda así:

\begin{lstlisting}
Linear
  |
  v
Issue
  |
  v
Branch pequeña
  |
  v
Desarrollo
  |
  v
Pull Request
  |
  +----> CI
  |
  +----> Code Review
  |
  v
Squash and Merge
  |
  v
main
  |
  v
Deploy

La branch se conserva para mantener visible el historial.
\end{lstlisting}

% -------------------------------------------------

\section{Referencia tecnológica}

Este tutorial se ocupa del flujo de trabajo. El stack, la arquitectura, la
estructura del monorepo, las variables de entorno y los despliegues están
documentados en \texttt{tech.tex}.

Consulten esa guía cuando una tarea necesite una decisión técnica. Si la
decisión afecta a ambos equipos, además de documentarla allí deben registrarla
en Linear y coordinarla con TI4 mediante el flujo descrito en este documento.

% -------------------------------------------------

\section{Principios generales}

Como TI4, esperamos que TI2 trabaje con estos principios:

\begin{enumerate}
  \item Registren el trabajo relevante para que podamos seguirlo.
  \item No hagan push directo a \texttt{main}.
  \item Mantengan las ramas pequeñas y de corta duración.
  \item Pidan una revisión antes de integrar el código.
  \item Aprovechen la revisión de TI4 para ganar autonomía, no para delegar su
    trabajo.
  \item Mantengan Linear actualizado para que refleje el estado real del
    trabajo.
  \item Usen GitHub para revisar y mantener el código.
  \item Mantengan todos los componentes principales del producto en el
    monorepo.
  \item Separen Presentación, Aplicación, Dominio e Infraestructura.
  \item Mantengan las reglas de negocio fuera de las aplicaciones cliente.
  \item Mantengan explícitos y versionados la API compartida y los paquetes
    internos.
  \item Avísennos pronto cuando encuentren un bloqueo o una decisión que no
    puedan resolver.
\end{enumerate}

% -------------------------------------------------

\section{Resumen}

Si tienen dudas durante el desarrollo, recuerden esta división:

\begin{center}
  \textbf{Linear = planificación y seguimiento}
\end{center}

\begin{center}
  \textbf{GitHub = código, revisión, CI e integración}
\end{center}

\begin{center}
  \textbf{Monorepo = código compartido, trazable y coordinado}
\end{center}

TI4 les entrega contexto, revisión y acompañamiento; TI2 implementa sus tareas
y gana autonomía con cada ciclo. El objetivo es que ambos equipos puedan
mantener el proyecto sin depender de una sola persona.

\end{document}
